\section{Zadanie 6}
Niech $X,Y$ będą niezależnymi zmiennymi losowymi o rozkładzie wykładniczym
z parametrem $1$. Obliczyć
\[
\lim_{R\to\infty} P\bigl(\min(X,Y)>R \mid X+Y>3R\bigr).
\]


\subsection*{Rozwiązanie}
\[
P(\min(X,Y) > R \mid X+Y > 3R)
=
\frac{P(\min(X,Y)>R \cap X+Y > 3R)}{P(X+Y>3R)}
\]

Policzmy prawdopodobieństwo z licznika:\\

\[
\begin{aligned}
P(\min(X,Y)>R \cap X+Y>3R)
&= \int_{R}^\infty f_X(x)
\int_{\max(3R-x,\,R)}^\infty f_Y(y)\,dy\,dx \\
&= \int_R^{2R} e^{-x}
\int_{3R-x}^\infty e^{-y}\,dy\,dx
+ \int_{2R}^\infty e^{-x}
\int_R^\infty e^{-y}\,dy\,dx \\
&= R e^{-3R} + e^{-3R}
\end{aligned}
\]

Policzmy prawdopodobieństwo z mianownika:\\

Zmienna $X+Y$ ma rozkład $\Gamma(2,1)$ (twierdzenie z wykładu) zatem:\\

\[
f_{X+Y}(x) =
\begin{cases} 
x e^{-x} & \text{dla } x \geq 0,\\
0 & \text{wpp.}
\end{cases}
\]

\[
\begin{aligned}
P(X+Y > 3R)
&= 1 - P(X+Y \leq 3R)
= 1 - \int_0^{3R} x e^{-x}\,dx \\
&= 1 - (-3R e^{-3R} - e^{-3R} + 1)
= e^{-3R} + 3R e^{-3R}
\end{aligned}
\]

\[
\begin{aligned}
\lim_{R\to \infty} P(\min(X,Y) > R \mid X+Y > 3R)
&= \lim_{R\to \infty}
\frac{e^{-3R}(1 + R)}{e^{-3R}(1 + 3R)}
= \frac{1}{3}
\end{aligned}
\]


\subsection*{Odpowiedź}
\[
\boxed{\frac{1}{3}}
\]
